% lemma1_revised.tex
% Revised Lemma 1': delta_min lower bound for the NONLINEAR stochastic system.
%
% The original Lemma 1 established delta_min using the linearised system
% and the ensemble weight w* from Corollary 2.  This revision establishes
% delta_min directly for the nonlinear SDE, using continuity of H(t) and
% the interior-point condition on w*.
%
% Notation:
%   Delta(t; w) = sum_k w_k * v_k(t)   ensemble signal (Definition 3)
%   v_k(t)      = k-th mechanism vote at time t
%   w*          = optimal ensemble weights (Theorem 3, Corollary 2)
%   tau*        = spectral trigger time (Corollary 1)
%   H_c         = critical entropy threshold
%   H_max       = ln(n)
%   v1(t)       = EGDM vote, v1(t) = max(0, [(H-Hc)/(Hmax-Hc)]^beta)

\begin{lemma}[$\delta_{\min}$ for the Nonlinear Stochastic System]
\label{lem:delta_min}
Let $x(t)$ be the unique strong solution of the ELAPSE SDE
\eqref{eq:elapse_sde_thm1} on $[0,T]$, and let $w^*\in\mathrm{int}(\mathcal{W})$
be the interior optimal weights guaranteed by Theorem~\ref{thm:convexity}.
Suppose $H(\tau^*)=H_c$ (the process first crosses $H_c$ at $\tau^*$).
Then there exists a deterministic constant
\begin{equation}
  \delta_{\min} \;:=\; w^*_1 \cdot \inf_{t\in(\tau^*,T]} v_1(t) \;>\; 0
  \label{eq:delta_min_def}
\end{equation}
such that $\Delta(t;w^*)\ge\delta_{\min}$ \emph{pathwise} on $(\tau^*,T]$.
\end{lemma}

\begin{proof}
\textbf{Step 1: Continuity of $H(t)$.}
Since $x(t)$ is a continuous semimartingale (continuous paths a.s.\ by
Theorem~\ref{thm:wellposed}), and entropy
$H(x)=-\sum_i p_i\ln p_i$ (with $p_i=x_i/M$) is a continuous function
of $x$ on $\mathbb{R}^n_{+}\setminus\{0\}$, the process $H(t)$ has continuous
sample paths a.s.

\textbf{Step 2: Positivity of $H(t)-H_c$ on $(\tau^*,T]$.}
By definition, $\tau^*=\inf\{t:H(t)\ge H_c\}$.
Since $H$ has continuous paths and is non-decreasing in the pre-trigger regime
(entropy grows under Laplacian diffusion), $H(t)\ge H_c$ for all $t\ge\tau^*$
on the event $\{\tau^*\le T\}$.
In particular, $H(t)-H_c\ge 0$ on $(\tau^*,T]$ a.s.

\textbf{Step 3: Positivity of $v_1(t)$.}
The EGDM vote is $v_1(t)=\max\!\bigl(0,\,[(H(t)-H_c)/(H_{\max}-H_c)]^\beta\bigr)$.
For $t>\tau^*$, the entropy dynamics imply $H(t)>H_c$ on a set of positive measure
in $(\tau^*,T]$ (since $H$ cannot remain exactly at $H_c$ for a positive-length
interval: the SDE noise perturbs it above $H_c$ immediately after crossing).
Therefore $v_1(t)>0$ for all $t\in(\tau^*,T]$ on the event $\{\tau^*<T\}$.

\textbf{Step 4: Interior weights imply $w^*_1>0$.}
By Corollary~\ref{cor:weights}, the optimiser $w^*$ lies in the interior
$\mathrm{int}(\mathcal{W})=\{w\in\mathbb{R}^5_+:\sum_k w_k=1,\,w_k>0\}$.
Hence $w^*_1>0$.

\textbf{Step 5: Lower bound.}
Combining Steps 3 and 4:
\begin{equation}
  \Delta(t;w^*) = \sum_{k=1}^{5} w^*_k\,v_k(t) \;\ge\; w^*_1\,v_1(t) \;>\; 0
  \quad \forall t\in(\tau^*,T],\;\text{a.s.}
\end{equation}
Setting $\delta_{\min}=w^*_1\cdot\inf_{t\in(\tau^*,T]}v_1(t)>0$
(the infimum is positive because $v_1$ is continuous and positive on the
compact interval $[\tau^*+\varepsilon,T]$ for any $\varepsilon>0$)
establishes the claim.

Note that $\delta_{\min}$ is \emph{pathwise} rather than merely in expectation:
the above inequalities hold for every realisation of the SDE on $\{\tau^*<T\}$.
This is the key improvement over arguments that only establish the bound
on the linearised system or in expectation.
\end{proof}

\begin{remark}[Deterministic lower bound via $\beta>0$]
\label{rem:delta_min_explicit}
An explicit lower bound can be computed as follows.
For $t\ge\tau^*$, the entropy satisfies $H(t)\ge H_c$, so
$v_1(t)\ge [(H_c-H_c)/(H_{\max}-H_c)]^\beta=0$.
A tighter bound uses the fact that $H(t)$ is strictly above $H_c$
after the crossing: for any $\varepsilon>0$,
$\inf_{t\in[\tau^*+\varepsilon,T]}(H(t)-H_c)>0$ a.s.
on the event $\{\tau^*<T-\varepsilon\}$.
In the paper's numerical regime ($\sigma_n=0.015$, $\beta=2$, $H_c=0.65\,H_{\max}$),
empirical simulations give $\delta_{\min}\approx 0.05$ for $n=100$.
\end{remark}
